\subsection{The wild local central groups and their actual lifting fibers}

\begin{statement}{E.74}{Proved}{Integral cohomology and every nonempty central fiber at three}
Retain the based cubic pro-three groups $W$ and their integral
graded lattices $\mathcal A_r(W)$ from E.41. Fix an embedding
of the algebraic closure of $\mathbb Q$ into an algebraic
closure of $\mathbb Q_3$, and put
\[
 k=\mathbb Q_3(\zeta_3),\quad
 \mathcal D=\operatorname{Gal}(k(3)/\mathbb Q_3),\quad
 D=\operatorname{Gal}(k(3)/k),\quad
 \Gamma=P\rtimes\langle c\rangle.
 \tag{E74.1}
\]
Here $k(3)$ is the maximal pro-three extension of $k$, and
$P=\langle\tau,\tau'\rangle$ is the actual global free
quotient from E.37, with $c\tau c=\tau'$. The localization
map $\pi_3:\mathcal D\to\Gamma$ is onto. It has a
continuous group section $j$ constructed below. Write
$R_3=\ker\pi_3$ and use the resulting actual coordinates
$\mathcal D=R_3\rtimes j(\Gamma)$.

For $A=\mathcal A_r(W)$ write $n_r=\operatorname{rank}A$,
$t_r=\operatorname{tr}J_r$, where $T_r,J_r$ are the
actual matrices of E.41. Define
\[
 \begin{split}
 \mathcal B_r(W)
   &=\operatorname{coker}\bigl(
       [\,T_r-7I\mid J_r+I\,]:A^2\longrightarrow A\bigr),\\
 b_1&=5,\qquad b_2=1,\qquad b_r=0\quad(r\ge3),\\
 d_r&=\frac{n_r+t_r}{2}+b_r.
 \end{split}
 \tag{E74.2}
\]
All cokernels have their compact quotient topology. There
are isomorphisms of compact $\mathbb Z_3$-modules
\[
 \begin{split}
 H^2_{\rm cts}(G_{\mathbb Q_3},A)&\simeq\mathcal B_r(W),\\
 H^1_{\rm cts}(G_{\mathbb Q_3},A)
   &\simeq\mathcal C_r(W)\oplus
       \operatorname{Hom}_{\Gamma,\rm cts}(R_3^{\rm ab},A)\\
   &\simeq\mathcal C_r(W)\oplus\mathbb Z_3^{d_r}.
 \end{split}
 \tag{E74.3}
\]
The first isomorphism is the local invariant pairing with
the invariant Tate dual, after a choice of compatible
three-power roots of unity. The second is restriction to
$j(\Gamma)$ and $R_3$; its inverse is the literal cocycle
\[
 a(rj(g))=\theta(r)+u(g),\qquad
 r\in R_3,\quad g\in\Gamma.
 \tag{E74.4}
\]
The last isomorphism uses a chosen basis of the actual
equivariant homomorphism lattice. Two normal generators
of $R_3$, explicit evaluation coordinates, and the inverse
coordinate map are specified in the proof. The resulting
coordinate choice is recorded separately from numerical
evaluation of a local relator.

For the first three degrees, all three geometric quotients
have the same answer. With $C_m=\mathbb Z/m\mathbb Z$,
the complete obstruction groups are
\[
 \begin{split}
 \mathcal B_1&=\mathbb Z_3^5,\\
 \mathcal B_2&=\mathbb Z_3\oplus C_3^{11}\oplus C_9^4,\\
 \mathcal B_3&=C_3^{63}\oplus C_9^{23}\oplus C_{27}^7.
 \end{split}
 \tag{E74.5}
\]
The complete central parameter groups are
\[
 \begin{split}
 H^1(G_{\mathbb Q_3},\mathcal A_1)
   &=\mathbb Z_3^{15}\oplus C_3^3,\\
 H^1(G_{\mathbb Q_3},\mathcal A_2)
   &=\mathbb Z_3^{46}\oplus C_3^9\oplus C_9^5,\\
 H^1(G_{\mathbb Q_3},\mathcal A_3)
   &=\mathbb Z_3^{330}\oplus C_3^7\oplus C_9^{56}
                         \oplus C_{27}^{25}\oplus C_{81}^3.
 \end{split}
 \tag{E74.6}
\]
In particular the additional free ranks in (E74.3) are
$9,22,164$. In degree four the full, marked and metabelian
groups have additional free ranks $1227,1222,741$,
respectively; their total $H^1$ ranks are $2475,2466,1485$.
Their torsion is precisely the corresponding E.41 torsion.

For every $r\ge3$, $\mathcal B_r(W)$ is finite and is killed by
\[
 3^{r+e_r},\qquad
 e_r=\max_{r-2\le j\le2r-2}\bigl(1+v_3(j)\bigr).
 \tag{E74.7}
\]
Thus the actual matrix in (E74.2) modulo
$3^{r+e_r+1}$ evaluates it by a finite Smith calculation.
The retained certificate gives the integer transformation
and inverse maps in degree one, and the modular transformation
and inverse maps, with proved sufficient precision, in
degrees two and three.

Let $W_r=W/\gamma_{r+1}W$. For each actual section
$\sigma:G_{\mathbb Q_3}\to W_r\rtimes G_{\mathbb Q_3}$,
write $\bar\sigma$ for its projection. The map
\[
 H^1(G_{\mathbb Q_3},\mathcal A_r(W))
   \longrightarrow
   \operatorname{Fib}_{[\bar\sigma]}
      \bigl(\operatorname{Sect}_{\mathbb Q_3}(W_r)/W_r
         \longrightarrow
         \operatorname{Sect}_{\mathbb Q_3}(W_{r-1})/W_{r-1}\bigr),
 \quad [z]\longmapsto[z\sigma]
 \tag{E74.8}
\]
is a homeomorphism of pointed spaces. Its inverse aligns
the lower sections by their unique geometric conjugator
and takes their central difference. Consequently (E74.3)
computes every nonempty central fiber at the wild place.
The obstruction of an arbitrary lower section takes values
in the explicitly computed group $\mathcal B_r(W)$.
Its individual nonlinear values remain part of E.23.

\emph{Manuscript: \S51.16, folios733v--736v, scans172--178,
formulas(371)--(379); Part E(c),(d) and follow-up(4).
The coefficient prime, geometric groups and arithmetic
source in this statement are exactly those displayed above.}
\end{statement}

\begin{proof}
\emph{The actual local source and an equivariant splitting.}
Restriction to the completion supplies the continuous map
$D\to P$. A proper closed image is contained in a subgroup
of index three: choose a finite three-group quotient in
which the image is proper, and then a maximal subgroup
of that finite group. The corresponding cyclic cubic
extension of $K=\mathbb Q(\zeta_3)$ is unramified outside
three and has trivial decomposition group at its unique
prime above three. It is therefore unramified at every
finite prime. The class number one computation in E.34
excludes this extension. This proves surjectivity.

The relative field $k(3)$ is stable under the automorphisms
of $k/\mathbb Q_3$, so (E74.1) gives an exact sequence
with quotient $C_2$. Its localization map onto $\Gamma$
is compatible with that quotient and is surjective by
the preceding argument. In the preimage of the specified
subgroup $\langle c\rangle\subset\Gamma$, choose a
complement to the normal pro-three kernel. The profinite
coprime complement theorem follows here by the finite
coprime complement theorem and compactness of compatible
finite complements. It gives an element $\widetilde c$
of order two mapping to $c$.

Choose any $\widetilde\tau\in D$ mapping to $\tau$ and
put $\widetilde\tau'=\widetilde c\widetilde\tau
\widetilde c$. The free property of $P$ extends these
two choices to a continuous group homomorphism $P\to D$.
Its composite with localization is the identity on the
two free generators and hence on $P$. Conjugation by
$\widetilde c$ interchanges the two chosen lifts, so this
homomorphism extends to the section $j:\Gamma\to\mathcal D$.
Every involution just used lies in the stated relative
quotient, with its map from the absolute local group.

By E.34 the based action on each $W$ factors through
$\Gamma$, and $R_3$ consequently acts trivially on $W$
and $A$. The local cyclotomic character also factors
through $\Gamma$: all three-power roots of unity occur
in the global extension unramified outside three.
In particular
\[
 \chi_3(j\tau)=7,\qquad \chi_3(jc)=-1.
 \tag{E74.9}
\]
These are the values of the specified lifted arithmetic
elements. All matrices below refer to those elements.

We prove the cohomological comparison with the full
absolute local group directly. The actual local/global
quotient is also recorded in the opened original
\cite[pp.~1--2; \S3]{E_WingbergFreeQuotients}.
Let $N_0=G_{k(3)}$. Kummer theory gives
$H^1(N_0,\mathbb F_3)=0$: every $a\in k(3)^\times$
lies in a finite Galois three-extension $E/k$.
Adjoining cube roots of all its $k$-conjugates gives
an elementary three-extension of $E$, since $\mu_3\subset E$.
The resulting splitting field is Galois of three-power
degree over $k$, so it already lies in $k(3)$.

Continuous cohomology of the closed subgroup $N_0$
is the direct limit over these finite fields $E$ of
$H^i(G_E,\mathbb F_3)$. This comparison is induced by
restriction: continuous finite-valued cochains extend
to a sufficiently small open neighborhood of the
closed subgroup, and their cocycle and coboundary
identities hold after shrinking that neighborhood.
For degree two, identifying $\mathbb F_3$ with $\mu_3$
identifies this system with $\operatorname{Br}(E)[3]$.
Its local invariant is multiplied by $[E':E]$ under
restriction to $E'$. A degree-three step of the
unramified tower beyond the finite residue degree
of $E$ gives $[E':E]=3$ inside $k(3)$, and kills
that group. Thus $H^2(N_0,\mathbb F_3)=0$.
All higher groups vanish by finite local duality
before passage to the same direct limit. These
are the restriction and invariant maps in
\cite[I, Theorem2.1 and Corollary2.3]{E_MilneADT}.

For every finite three-primary $D$-module $M$, the
$N_0$ action is trivial. A filtration of its underlying
finite abelian group with trivial $\mathbb F_3$
factors therefore proves $H^i(N_0,M)=0$ for $i>0$.
The Hochschild--Serre spectral sequence now gives
inflation isomorphisms $H^i(D,M)\to H^i(G_k,M)$
in every degree, with their actual restriction maps.
Exact $C_2$ invariants give the comparison over
$\mathbb Q_3$. The four Kummer classes below give
$\dim H^1(D,\mathbb F_3)=4$, so the pro-three
Frattini criterion supplies four topological
local generators. Local duality and the proved
comparison also recover its dimension-two
Demu\v skin pairing directly.
At the coefficient levels $A/3^mA$ all local cohomology
groups are finite, by local duality. Their inverse systems
are Mittag--Leffler, so every relevant derived-limit term
vanishes. Passing to the inverse limit gives the
continuous compact-coefficient comparisons
\[
 H^i_{\rm cts}(G_{\mathbb Q_3},A)
       \xleftarrow[\rm inflation]{\ \sim\ }
 H^i_{\rm cts}(\mathcal D,A),\qquad i=1,2,
 \tag{E74.10}
\]
and vanishing above degree two. For degree one and for
the nonabelian section sets, factorization is also direct:
over $k$ the map $g\mapsto(a(g),\rho(g))$ has a
pro-three semidirect target. It kills $N_0$, so every
such cocycle factors through $\mathcal D$. Inflation and
factorization are inverse on cocycles and conjugations.

\emph{The integral degree-one splitting and its coordinates.}
For an additive cocycle $a$ on $\mathcal D$, restriction
to $R_3$ is a continuous homomorphism $\theta$. The
cocycle identity and the trivial $R_3$ action give
\[
 \theta(j(g)rj(g)^{-1})=g\theta(r).
 \tag{E74.11}
\]
Restriction to $j(\Gamma)$ is a cocycle $u$. Conversely,
the multiplication law of $R_3\rtimes j(\Gamma)$ and
(E74.11) verify (E74.4). A coboundary vanishes on
$R_3$ and restricts to the ordinary coboundary on
$\Gamma$. Thus these restriction maps and (E74.4)
give the middle isomorphism in (E74.3), integrally.
E.41 computes the first factor by normalizing the
value at $c$ to zero and evaluating at $\tau$:
\[
 H^1(\Gamma,A)=A/(T_r-I)A^c=\mathcal C_r(W),
 \qquad A^c=(I+J_r)A.
 \tag{E74.12}
\]
Division by two proves the last equality.

Here are actual coordinates for the second factor.
Set $\lambda=1-\zeta_3$, $\mathcal O=\mathbb Z_3[\zeta_3]$
and $U_i=1+\lambda^i\mathcal O$. The unit filtration gives
$U_1=\mu_3\times U_2$: the quotient $U_1/U_2$ has order
three, and $\zeta_3\notin U_2$. The logarithm and
exponential converge inversely on $U_2$ and
$\lambda^2\mathcal O$, since $v_3(\lambda^2)=1>1/2$.
Their valuations follow term by term from the factorial
and integer denominators. Put $u_2=1+\lambda^2$,
$u_3=1+\lambda^3$. Their logarithms are respectively
$\lambda^2,\lambda^3$ modulo $\lambda^4\mathcal O$.
For example the denominator-three term of
$\log(1+\lambda^2)$ already has $\lambda$-valuation
four; every later term has at least that valuation.
Since $3\lambda^2\mathcal O=\lambda^4\mathcal O$ up
to a unit, the four Kummer classes
\[
 \lambda,\quad\zeta_3,\quad u_2,\quad u_3
 \quad\text{form a basis of }k^\times/k^{\times3}.
 \tag{E74.13}
\]
The first two span precisely the restrictions of
$H^1(P,\mathbb F_3)$, by the explicit global S-unit
calculation of E.34. The split five-term sequence and
$H^2(P,\mathbb F_3)=0$ show that the restrictions of
$\kappa_{u_2},\kappa_{u_3}$ are a basis of
$H^1(R_3,\mathbb F_3)^P$. Choose elements $r_1,r_2\in R_3$
dual to that basis. They normally generate $R_3$ in $D$.
Indeed, a nontrivial quotient by their closed normal
closure would have a nontrivial finite normal three-group
image. Its Frattini quotient has nonzero coinvariants
under the finite acting three-group, and hence a nonzero
invariant character. That character vanishes on $r_1,r_2$,
contrary to their chosen dual coordinates.

Let $\Lambda=\mathbb Z_3[[P]]$. Conjugation through $j$
and the two generators give an actual continuous module
surjection $\Lambda^2\to R_3^{\rm ab}$. Evaluation
therefore embeds the second factor of (E74.3) as
\[
 E:\operatorname{Hom}_{\Gamma,\rm cts}(R_3^{\rm ab},A)
       \lhook\joinrel\longrightarrow A^2,
 \qquad E(\theta)=(\theta(r_1),\theta(r_2)).
 \tag{E74.14}
\]
Its image $\mathcal E_A$ consists of pairs whose
$\Lambda$-linear map kills the actual relation submodule
and respects the actual $c$ action. These are closed
integral linear equations. The image is closed and
saturated: when three times a pair satisfies all those
equations, torsion-freeness of $A$ implies that the pair
satisfies them. Thus $\mathcal E_A$ is a finite free
direct summand of $A^2$.

Choose a unimodular map $S_A:A^2\to A^2$ carrying this
actual submodule to the first $d_r$ coordinates, whose
rank is proved below. The extra coordinates of $[a]$
are the first $d_r$ entries of
$S_A(a(r_1),a(r_2))$. Conversely append zeros, apply
$S_A^{-1}$, and take the unique equivariant homomorphism
having that pair of values. Formula (E74.4) gives the
inverse class together with its E.41 coordinate.
All choices and both maps have now been specified.
The numerical entries of $S_A$ depend on the actual
relation submodule; these entries are not among the
Smith computations in (E74.5). The evaluated module
type and rank below require no unevaluated relator value.

The topological assertions follow from the same maps.
Cocycles on the finitely generated pro-three group $D$
form a closed submodule of the finite product of copies
of $A$ obtained by evaluating on its generators.
The extra $C_2$ equations are closed. Coboundaries are
compact images of $A$. Thus all the degree-one groups
are compact Hausdorff, and the displayed continuous
bijections are homeomorphisms.

\emph{Local duality, with its exact twist.}
At each finite level the cup product followed by the
local Brauer invariant gives the perfect duality
\[
 H^2(G_{\mathbb Q_3},A/3^mA)\times
 H^0(G_{\mathbb Q_3},\operatorname{Hom}(A/3^mA,\mu_{3^m}))
       \longrightarrow\tfrac1{3^m}\mathbb Z/\mathbb Z.
 \tag{E74.15}
\]
This is the specialization of
\cite[I, Theorem2.1 and Corollary2.3, pp.~26--28]{E_MilneADT}.
The maps under coefficient reduction are dual to
the inclusions of the Tate-dual coefficient groups.
Their inverse limit consequently identifies $H^2$ with
the compact coinvariants of $A(-1)$. Explicitly, a
coinvariant is paired with an invariant element of
$A^*(1)\otimes\mathbb Q_3/\mathbb Z_3$ by evaluation;
the inverse is the unique cohomology class giving that
functional under (E74.15). This defines both directions
of the first isomorphism in (E74.3).

The localization map is onto $\Gamma$, so these
coinvariants are precisely the $\Gamma$ coinvariants.
On $A(-1)$ the two topological generators $\tau,c$
act as $7^{-1}T_r,-J_r$. Their relation submodule is
therefore $(T_r-7I)A+(J_r+I)A$. Relations for products
and inverses follow from those for the generators;
continuity gives the relations for their closure.
This image is already compact and closed. This proves
the exact matrix formula (E74.2).

Its rational free rank is $b_r$. To compute it, use
the actual split rational lattice from E.40--E.41.
On degree one the Frobenius at the $K$-prime above two
acts as $-2$ on the elliptic part and as $4$ on the
toric part. The toric character module is the
degree-zero cusp module on nine cusps; $c$ has three
fixed cusps and three pairs. After the Tate twist its
coinvariant dimension is $6-1=5$. The elliptic part
has no eigenvalue four, proving $b_1=5$. In degree two
the only four-eigenspace is
$\bigwedge^2T_3E\otimes\mathbb Q_3=\mathbb Q_3(1)$,
with its actual Weil-pairing action. It contributes
one dimension, proving $b_2=1$. All degrees at least
three have eigenvalues $(-2)^j$, $r\le j\le2r$;
none is four. The exact graded quotients of E.41
retain these weight spaces and hence have $b_r=0$.

\emph{The integral rank and all degree-one torsion.}
The element $j(\tau)$ acts by the positive-weight
Frobenius matrix $T_r$, which has no eigenvalue one.
Thus $H^0(G_{\mathbb Q_3},A)=0$. Write $R_A$ for the
free rank of $H^1(G_{\mathbb Q_3},A)$; finite generation
was proved in (E74.12)--(E74.14). For $m$ beyond the
exponents of its torsion and of $\mathcal B_r$, the
coefficient exact sequence gives
\[
 \begin{split}
 H^0(G_{\mathbb Q_3},A/3^m)&=H^1(G_{\mathbb Q_3},A)[3^m],\\
 0\longrightarrow H^1(A)/3^m&\longrightarrow H^1(A/3^m)
       \longrightarrow H^2(A)[3^m]\longrightarrow0,\\
 H^2(G_{\mathbb Q_3},A/3^m)&=H^2(G_{\mathbb Q_3},A)/3^m.
 \end{split}
 \tag{E74.16}
\]
In the middle row the arithmetic group is the same
$G_{\mathbb Q_3}$ throughout. The finite local Euler
characteristic theorem
\cite[I, Theorem2.8, pp.~31--34]{E_MilneADT} says that
the alternating product of these orders is $3^{-mn_r}$.
The two finite torsion orders cancel in (E74.16), giving
$3^{m(b_r-R_A)}=3^{-mn_r}$. Therefore
\[
 R_A=n_r+b_r,\qquad
 \operatorname{rank}\mathcal E_A
   =n_r+b_r-\frac{n_r-t_r}{2}=d_r.
 \tag{E74.17}
\]
The direct summand $\mathcal E_A$ is torsion-free, so
every torsion summand of local $H^1$ is exactly the
E.41 torsion. This proves (E74.6) and all stated
degree-four ranks.

\emph{Finite precision and the retained Smith maps.}
Let $F$ denote the actual Frobenius at the $K$-prime
above two. The split lattice of E.41 has maps $B_r,C_r$
with $B_rC_r=3^rI$ and
$F B_r=B_r F_{\rm split}$. Its diagonal eigenvalues
are $(-2)^j$, $r\le j\le2r$. For $r\ge3$,
\[
 v_3((-2)^j-4)=1+v_3(j-2).
 \tag{E74.18}
\]
The scalar four is a three-adic unit. It follows that
$3^{e_r}(F_{\rm split}-4I)^{-1}$ is integral. Multiplying
by $B_r,C_r$ shows that
$3^{r+e_r}(F-4I)^{-1}A\subset A$. The image $(F-4I)A$
is zero in the coinvariants (E74.2), because $\chi_3(F)=4$.
This proves the annihilator bound (E74.7), including its
passage through each actual geometric quotient.

For degree two the free coordinate is the first Hall
coordinate $[x_1,x_0]$, the projection to
$\bigwedge^2T_3E$. The entire first row of the matrix
in (E74.2) is zero. On its coordinate-zero lattice
$K_2=\langle e_1,\ldots,e_{44}\rangle$, the supplied
integer matrices verify
\[
 (F-4I)B_2D_2C_2\big|_{K_2}=108I_{K_2},
 \qquad
 D_2\big|_{\text{eigenvalue }e}
   =\begin{cases}0&e=4,\\12/(e-4)&e=-8,16.
     \end{cases}
 \tag{E74.19}
\]
The elliptic coordinate of $C_2K_2$ is zero, which
verifies the identity without a rational projection
error. Hence the torsion is killed by $3^3$, and
precision $3^4$ suffices.

The program \texttt{e\_wild\_layers.py} constructs the
actual matrices by integral Hall substitution from
E.41. It checks equality with both quotient matrices
in degrees one through three. In degree one it retains
an integer Smith decomposition and both inverse
transformations. In degrees two and three it retains
every elementary operation, the left and right
transformations, both inverses, and the exact modular
Smith identities at moduli $81$ and $729$. The
nonzero Smith valuations and free ranks are
\[
 \begin{array}{c|r|rrrr|r}
 r&n_r&v=0&v=1&v=2&v=3&\text{free rank}\\\hline
 1&10&5&0&0&0&5\\
 2&45&29&11&4&0&1\\
 3&330&237&63&23&7&0.
 \end{array}
 \tag{E74.20}
\]
These are (E74.5). The annihilator bounds exclude
undetected elementary divisors beyond the precision.
For a finite factor, multiply by the stored left
transformation and reduce each coordinate modulo its
displayed power of three. For the inverse, lift those
digits, apply the stored inverse, and pass to the
cokernel. Differences between lifts are killed by
the proved annihilator. In degree two first subtract
the exact free coordinate times $e_0$ and apply this
construction to $K_2$; retain the first coordinate
in $\mathbb Z_3$. This supplies the actual coordinate
maps, including the free part.

\emph{Section fibers, stabilizers, and the comparison maps.}
A geometric element fixed by a section-twisted action
on $W_{r-1}$ has a first nonzero graded coordinate.
Inner conjugation by the section's geometric part
acts trivially on that grade, so this coordinate is
fixed by $T_i$. Its absence of eigenvalue one is a
contradiction. Thus every such stabilizer is trivial.

Fix the actual section $\sigma$ in (E74.8). Multiplying
it by an $A$-cocycle gives a section because $A$ is
central in $W_r$ and has the original arithmetic
action. Another section whose lower class is
$[\bar\sigma]$ has a unique lower geometric
conjugator aligning the two lower sections. Lift
that conjugator to $W_r$ and align. Their difference
is then an $A$-cocycle. Changing the lifted conjugator
changes this cocycle by an $A$-coboundary. Conversely
two central differences yielding conjugate sections
have a conjugator projecting to a lower stabilizer;
that projection is the identity, so their difference
is an $A$-coboundary. This proves both directions
and injectivity of (E74.8). All section spaces factor
through the finitely generated relative source
$\mathcal D$ and are compact; their geometric orbit
spaces are Hausdorff. The continuous bijection in
(E74.8) is consequently a homeomorphism.

For completeness, the precise obstruction map of a
lower geometric cocycle $a$ is obtained with the
continuous Hall lift $h_r:W_{r-1}\to W_r$ of E.41:
\[
 (g,h)\longmapsto
 h_r(a(g))\,g(h_r(a(h)))\,h_r(a(gh))^{-1}\in A.
 \tag{E74.21}
\]
Associativity gives its additive two-cocycle identity;
changing the Hall lift changes it by a coboundary.
Its local invariant functional (E74.15), followed
by the proved Smith map, is its value in $\mathcal B_r$.
Vanishing means that an actual central correction
of $h_r\circ a$ is a cocycle, with inverse provided
by the same identity. This specifies the obstruction
and its map while leaving its unevaluated nonlinear
values in the outstanding classification.

Finally the actual epimorphisms from the full group
to its metabelian and marked quotients induce the
graded maps $q_r:A\to A'$. These commute with
$T_r,J_r$ and with the local source. On (E74.3) the
map is $([u],\theta)\mapsto([q_ru],q_r\theta)$;
on (E74.2) it is the map of cokernels induced by
$q_r$. Formula (E74.21) commutes with these maps
at the cohomology-class level, since changing a
chosen lift only changes a coboundary. Thus the
computed wild layers retain the original geometric
comparison maps and every stated ambiguity.
\end{proof}
